\section{Evaluation Plan}
\label{sec:evaluation}

This section is a protocol for future experiments. It contains no generated
measurements or assumed speedups. The companion experiment matrix defines
E00--E16, including instruments, negative cases, and execution gates.
Table~\ref{tab:rqs} connects the principal questions to those experiments.
The existing passthrough observations in Section~\ref{sec:agent-observations}
supply a measured request-structure pilot for RQ1; they do not complete
E01's locality/sharing study or any system-performance comparison.

\begin{table*}[t]
\centering
\scriptsize
\begin{tabularx}{\textwidth}{@{}p{0.65cm}p{4.6cm}Yp{1.60cm}@{}}
\toprule
RQ & Question & Primary evidence & Plan\\
\midrule
1 & Which workload properties favor shared immutable state? & Observed: per-task FUSE request mix, read/write sizes, and time variation. To measure: cross-task read-set overlap, mutation footprint, locality, and tool-phase-labelled access. & E01--E02\\
2 & When do delta-based workspaces reduce workspace cost? & Time to usable workspace, private metadata, modification and lifecycle costs & E03--E04\\
3 & What does packed metadata save, and at what cost? & Metadata-KV calls, object requests/bytes, latency, cache effects, construction amortization & E05--E07\\
4 & What changes with $g$, $(B,q)$, and client read policy? & Object formation, persisted geometry, physical bytes, verification/EC CPU, durable writes and repair & E08--E10\\
5 & Do the mechanisms provide additional combined value? & Legal factorial or nested comparisons under one runtime contract & E11\\
6 & Do the effects matter for complete multi-agent tasks? & Deterministic replay, live-task makespan, resource use, output-quality controls & E12--E13\\
7 & Which semantics and failure boundaries hold? & Differential operations, fault injection, integrity, reclamation and long-run resource bounds & E14--E16\\
\bottomrule
\end{tabularx}
\caption{Research questions and evidence status. RQ1 separates the observed
request-structure pilot from further measurements; the remaining rows are
planned. E00 first checks that each
configuration is connected and produces the claimed effective treatment.}
\label{tab:rqs}
\end{table*}

\subsection{Environment and baseline contracts}

Our proposed primary testbed has two client machines, one dedicated metadata
machine, and four RustFS storage machines connected by 25 GbE. A practical
target is 24--32 physical cores and 128 GiB memory per client, with independent
NVMe storage and fixed cache allocations. The initial storage geometry is
one drive per storage host in a verified 2+2 set. These are requested resources,
not a report of machines already used. A single-host four-drive setup serves
only as a functional and mechanism testbed; it cannot establish independence
of machine failure domains.

We will record processor, memory, disks, network, OS/kernel, FUSE settings,
toolchain, complete configuration, build features, immutable image digests,
and source revisions. Client object-block granularity $g$, durability,
compression, cache budget, readahead,
direct I/O, metadata TTL, inline thresholds, and backend gates are part of
the treatment definition. Larger coding geometries will be separate
experiments because they can change whether the selected $q$ is smaller than
the logical shard size.

External baselines include full copying, reflinks where supported, OverlayFS,
BranchFS, and an object-backed file system such as JuiceFS~\cite{juicefs_architecture} sharing the same
backend resources. At least one mature immutable lower, such as Nydus or
CernVM-FS, will be tested with a compatible writable layer. An immutable
sorted-KV lower with equal records and cache capacity provides a specific
layout control. A real DeltaFS comparison depends on its operation and
workflow compatibility; we will not label our sorted-KV control ``DeltaFS.''

We will compare local workspace service separately from remote durable
service. An OverlayFS local fsync is not equivalent to remote coded-object
completion. Similarly, BranchFS's ephemeral branch uses different fsync and
commit semantics~\cite{branchfs2026}. Where a baseline needs a checkpoint to
meet the selected persistence boundary, its completion cost is included.
If equivalent semantics cannot be provided, we report separate operating
points instead of one unconditional speed ranking.

\subsection{Workloads and measurements}

\researchq{1} The workload study will collect repository navigation, reading,
editing, build/test, and cleanup activity from a pinned agent harness and
model configuration. It will measure read-set overlap, directory locality,
reuse distance, mutation footprint, and per-task operation distributions.
Expected lookup and attribute probes are distinguished from I/O failures.
An instrumentation-on/off control will quantify trace overhead and loss.
The pilot motivates retaining Core and Exploratory task strata, both
successful and negative metadata requests, mixed request sizes, and temporal
bursts. Its aggregate method counts and size histograms are insufficient for
semantics-preserving replay: the next capture must retain paths, payloads,
handle dependencies, synchronization relations, and tool-phase boundaries.
It must also connect FUSE request sizes to client object-block granularity
$g$, final uploaded object lengths $L_o$, and actual client fetches before
testing $(B,q)$ policies. These observations distinguish object formation
from request selection and backend processing.

The proposed expanded panel contains 20 compatible tasks each from
SWE-bench Verified, a fixed public SWE-bench Pro version~\cite{swebench_pro_v2}, and a multilingual
set informed by Multi-SWE-bench~\cite{swebench2024,multi_swebench2025}.
Exact task lists and revisions will be frozen after capability screening.
An official audit has identified limitations in using SWE-bench Verified
to measure frontier model ability~\cite{swebench_audit2026}; we use it as a
reproducible workload source and retain other task families. Filesystem
performance does not establish improved reasoning ability.

The minimum live-agent core will use 20 prespecified tasks, three systems,
and three independent repetitions at one concurrency setting: 180 runs.
The 60-task extension requires 540 runs for that same setting, before any
additional model or concurrency sweep. Pilot task duration and token use
will determine a frozen budget. Broad concurrency sweeps first use replay;
they are not silently multiplied across every live condition.

To make the file-system load explicit, each task is segmented into six
traceable phases: (A) workspace creation and dependency discovery, (B)
repository navigation (`stat', `readdir', `open', `readlink', and small
reads), (C) edit and patch writes, (D) build and test with sequential reads
and generated files, (E) dependency installation or cleanup with create,
rename, unlink, and hard-link operations, and (F) diff, discard, or retained
workspace teardown. We retain the phase boundaries in the trace and report
both per-phase and whole-task critical paths. A controlled replay uses the
same paths, payload hashes, handle lifetimes, and synchronization events;
the live run additionally records model/tool spans. This prevents a fast
provisioning path from being mistaken for a faster build, and exposes the
small-file metadata pressure that a fio-only study would miss.

The primary agent-load panel is a small, prespecified matrix rather than a
post-hoc selection of favorable traces:

\begin{table*}[t]
\centering
\scriptsize
\begin{tabularx}{\textwidth}{@{}p{1.1cm}p{3.3cm}Yp{2.0cm}@{}}
\toprule
Phase & Representative operations & Treatment and baseline controls & Primary outputs\\
\midrule
A/B & fan-out, path walk, `git status', grep, metadata-only scans &
shared lower versus materialized checkout; cold/index-warm/data-warm states &
ready time, lower-KV calls, GETs, metadata P99\\
C & patch, overwrite, truncate, sparse write, fsync &
matched write-completion and cache budgets; 0.1--10\% mutation &
upper bytes, commit and drain tail, write amplification\\
D & compiler reads, test binaries, generated artifacts &
same dependency image and local build cache; sequential/random mix &
critical path, read bytes, CPU, task makespan\\
E/F & install, rename/unlink storms, retry, diff/discard &
same lifecycle and durability boundary; retain failed and timed-out runs &
whiteouts, retained objects, cleanup lag, success\\
\bottomrule
\end{tabularx}
\caption{Agent file-system load panel. The table defines the workload
phases; it does not assert that every agent task exercises them equally.}
\label{tab:agent-load}
\end{table*}

\researchq{2} Provisioning experiments vary lower size, workspace count, and
modification fraction. Initial values are $10^4$--$10^6$ files, 1--128
workspaces, and 0--50\% modified files. We measure control-plane return,
mount readiness, first navigation, and first durable edit separately.
Prepositioned and cold-staged baselines are separate cases. Churn experiments
add deletion, rename, truncation, and repeated lifecycles, subject to the
selected implementation's ownership and cleanup support.

\researchq{3} Metadata lookup experiments compare mutable per-file KV,
immutable sorted-KV, and the packed layout using equivalent values and
cache budgets. Standalone read-only access tests the packed representation
independently of delta-based workspaces; workspace experiments then compare
the same representations as alternative lowers. Both include positive and
negative lookup, direct getattr, enumeration,
lookup-plus-read, deep paths, and wide directories. The snapshot reader's
large-directory study is separated from the Entity bridge's 4,096-entry
limit. We count upper/control/lower KV requests independently, as well as
index, GroupMeta, descriptor, and payload object requests and bytes.

Cold-client, index-warm, metadata-warm, and data-warm states are explicitly
defined; backend page-cache state is another variable. Unsupported
payload-cache controls are not treated as available flags. Snapshot
construction, publication, rebuilding, and reclamation enter each resource
accounting metric. We report a break-even reuse count only if the measured
net benefit per reuse is positive, with the reuse unit fixed for each
standalone or workspace experiment.

\researchq{4} A backend-only study uses a fixed S3 request stream before adding
the file-system path. Object sizes cross inline, 1 MiB, encoding-block, and
multipart boundaries. $B$ settings include default, the current automatic hint
policy, and the four supported fixed values. $q$ settings are legacy or the implemented
selector, grouped by \emph{persisted} geometry. If automatic and default
produce the same $(B,q)$, they are a no-treatment control, not evidence for a
benefit from the automatic hint policy.

The file-system study adds $g$ as a separate object-formation factor. A
controlled sweep holds the logical write stream constant and records final
object lengths $L_o$, object counts, and uploaded bytes, including compression,
headers, and tails. It first holds backend geometry fixed, so changes caused
by forming different objects are not attributed to $B$ or $q$. Subsequent
comparisons cross these factors only where the selected upload path supports
them; a fixed configuration does not establish an automatic selection policy.
The ordinary client's range-read threshold scales with $g$. A causal sweep
must hold the absolute threshold constant where supported, or expose its
change as a coupled client-policy factor.

Reads vary length, alignment, crossing of encoding-block and $q$-interval
boundaries, and healthy versus degraded state. The primary measurements are physical data/parity
read bytes, checksum work, EC CPU, HTTP payload, request counts, and latency.
Client range-read and whole-object fetch policies are then separated while keeping
cache admission and coalescing constant or exposing them as additional
factors. Upper and lower results are separate because their integrity
and upload paths differ.

Writes include append, overwrite, tails, multipart retry, and frequent
durability requests. We record accepted writes, staging, remote completion,
metadata commit, and drained background debt. Repair experiments measure
read/write/network work and interference with healthy requests. The paper
will claim lower write amplification only if these measurements support
it; a read benefit can coexist with increased checksum or write cost.

\subsection{Composition and complete agent execution}

\researchq{5} The mechanisms address separate choices: workspace view
composition, frozen metadata representation, and object formation and layout.
Their comparisons must retain one runtime contract. Switching
from a buffered ordinary volume to a direct-I/O Entity mount is not a pure
delta-based workspace ablation. We will implement or validate matched controls for complete
metadata materialization versus shared lower state, metadata representation,
$g$, and backend geometry, with client read policy held constant or varied
explicitly. Packed metadata supports standalone read-only use and can also
serve as an optional workspace lower; neither role makes delta-based
workspaces a prerequisite for packing. Only configurations that can actually
be constructed enter a factorial comparison. If a selected entry point wires
a treatment only within a particular workspace path, its comparison is
nested within that runtime because of the implementation connection.
Where full crossing is possible, task-level log-latency interactions will
test whether combined effects differ from independent component effects.

\researchq{6} Deterministic replay and live agents answer different questions.
Replay freezes tool inputs, write contents, and dependencies, preserving
file-tree and output checks. A FUSE log containing only inode IDs and
request lengths is not necessarily replayable; paths, handle lifetimes,
payload references, initial state, and ordering must be recovered or
recaptured. Arrival-driven and closed-loop replay are reported separately.

Live AX/DSH experiments keep the task, model identifier, prompt, tools,
token/time budgets, and concurrency policy fixed. Mount probes must verify
that workspaces and subprocesses reach the intended file system rather than
a fallback directory. We measure batch makespan, tasks per hour,
provisioning, resource use, and the traced storage critical path. Model
waiting and I/O can overlap, so their durations cannot simply be added.
Output correctness is checked with the task harness; an exit code of zero
is not a resolved issue.

\subsection{Correctness, failures, and analysis}

\researchq{7} Differential and known-answer tests cover whiteouts, same-name
recreation, lower-file modifications, truncate/extend, holes, links,
rename-overwrite, open-unlink-read, and cross-workspace visibility within
the supported contract. Fault injection targets object upload, manifest
and binding publication, mapping commit, process restart, missing objects,
frame corruption, and repair interruption. Fresh-only Entity behavior is
reported as such; reopen or crash-takeover experiments require a completed
runtime extension. Fail-closed behavior is distinct from successful recovery.

Long runs cover multiple cache and cleanup cycles and track tail latency,
memory, open descriptors, owners, pending uploads, retained objects, and
reclamation backlog. Background publishing or repair is added only where
the selected runtime supports it. High mutation, poor locality, frequent
snapshot rebuilding, and large sequential reads are explicit boundary
workloads, not exclusions from the evaluation.

The independent unit is a run or a task, not each of its millions of
requests. Initial plans use at least ten independent microbenchmark blocks
and three live repeats per selected task; the formal sample sizes are
frozen after a variance pilot. Conditions are interleaved in randomized
host/time blocks. We report absolute effects, paired ratios, and 95\%
task- or repository-clustered bootstrap intervals. Repeated task outcomes
are not treated as independent binary pairs.

For live speed comparisons, the primary endpoint includes all assigned
tasks under the same deadline; failed and timed-out attempts remain in the
throughput and completion accounting. Completed-only latency is secondary
and accompanied by success rates and failure causes. Early failures do not
count as fast completions. This protocol prevents storage behavior or model
failures from creating an apparent speedup through selective omission.
